@@M0@@

Finally, let us indicate the origin of formulas (2.3.1) and (2.3.3). We consider Example 2.2.b), where the operations @@M1@@ are defined by the exterior powers of @@M2@@-Modules. We want @@M3@@ to be a @@M4@@-homomorphism; we must therefore define laws on @@M5@@ such that for all @@M6@@-Modules @@M7@@, we have:
@@M8@@
now relations (2.3.1) and (2.3.3) are precisely the preceding relations applied to the case where @@M9@@ and @@M10@@ are invertible modules.

We shall see in Exp. VI to what extent the desired property is satisfied.

\sgaStableTarget{sga6.e05.def.2-4}\textbf{Definition 2.4.} \uline{A pre}-@@M11@@-\uline{ring} @@M12@@ \uline{will be called a} @@M13@@-\uline{ring if the homomorphism of abelian groups} @@M14@@ \uline{is a} @@M15@@-\uline{homomorphism} \uline{for the pre}-@@M16@@-\uline{ring structure on the second member defined in} 2.3.

This therefore means that we have the relations: @@M17@@,
\sgaStableTarget{sga6.e05.eq.2-4-1}@@M18@@

In terms of the maps @@M19@@, the last two relations are equivalent to: @@M20@@, @@M21@@,
\sgaStableTarget{sga6.e05.eq.2-4-2}\sgaStableTarget{sga6.e05.eq.2-4-3}@@M22@@

